a) Der Großkreis ist

\mathdisp {G= \left\{ (x,y,z) \in S  \mid  ax+by+cz = 0        \right\}} {  }

und die beiden offenen Halbsphären sind

\mathdisp {U_+ = \left\{ (x,y,z) \in S  \mid  ax+by+cz > 0        \right\}} {  }

bzw.

\mathdisp {U_- = \left\{ (x,y,z) \in S  \mid  ax+by+cz < 0        \right\}} { . }

b) Für jede offene Halbkugel $U$ und den zugehörigen Großkreis $G$ ist
\mathl{U \cap G = \emptyset}{.} Der maximale Abstand von zwei Punkten
\mathl{P,Q \in S}{} ist
\mathl{2}{,} und dies ist genau dann der Fall, wenn die beiden Punkte gegenüber
\zusatzklammer {antipodal} {} {} liegen, wenn also ihre Verbindungsgerade durch den Kugelmittelpunkt geht
\zusatzklammer {also bei
\mathl{Q=-P}{}} {} {}. Ein solches antipodale Paar liegt nicht auf einer offenen Halbsphäre, da bei
\mathl{P=(x,y,z)}{} und
\mathl{P \in U_+}{} ja
\mathl{ax+by+cz>0}{} und daher
\mathl{a(-x) +b(-y)+c(-z) <0}{} gilt, also
\mathl{-P \in U_-}{.}

Wir nehmen nun an, dass es eine offene Überdeckung

\mathdisp {S \subseteq U_1 \cup U_2 \cup U_3} {  }

mit drei offenen Halbsphären
\mathl{U_1,U_2,U_3}{} gibt
\zusatzklammer {die entsprechenden Ebenen und Großkreise seien mit
\mathkor {} {E_i} {bzw.} {G_i} {} bezeichnet} {} {.}
Wegen
\mathl{U_1 \cap G_1 = \emptyset}{} folgt

\mathdisp {G_1 \subseteq U_2 \cup U_3} { . }

Der Durchschnitt
\mathl{G_1 \cap E_2}{} enthält mindestens zwei antipodale Punkte
\mathkor {} {P} {und} {-P} {}
Dabei ist
\mathl{P,-P \not\in U_2}{.} Da $U_3$ nach der Vorüberlegung kein antipodales Punktepaar enthält, gehört einer dieser Punkte auch nicht zu $U_3$ und wir haben einen Widerspruch.

 [[Kategorie:Latexseite]]
